\documentclass[12pt,a4paper]{article}
\usepackage[utf8]{inputenc}   
\usepackage[T1]{fontenc}      
\usepackage[spanish, es-tabla]{babel} 
\usepackage{geometry}
\geometry{top=2.5cm, bottom=2.5cm, left=3cm, right=3cm}
\usepackage{listings}
\usepackage{xcolor}
\usepackage{hyperref}

\definecolor{codegray}{rgb}{0.5,0.5,0.5}
\definecolor{backcolour}{rgb}{0.95,0.95,0.92}

\lstset{
    inputencoding=utf8,
    extendedchars=true,
    literate={á}{{\'a}}1 {é}{{\'e}}1 {í}{{\'i}}1 {ó}{{\'o}}1 {ú}{{\'u}}1 {ñ}{{\~n}}1 {Ñ}{{\~N}}1,
    backgroundcolor=\color{backcolour},
    basicstyle=\ttfamily\small,
    breaklines=true,
    keywordstyle=\color{blue},
    commentstyle=\color{codegray},
    stringstyle=\color{orange},
    keepspaces=true,
    showstringspaces=false,
    tabsize=2
}

\title{\textbf{Manual para Linux}}
\author{Gabriel Sandi}
\date{Febrero 2026}

\begin{document}

\maketitle
\tableofcontents
\newpage

\section{Objetivo}
Este manual técnico tiene como propósito servir de guía integral para el aprendizaje y consulta rápida sobre la gestión del sistema operativo Linux, cubriendo desde la navegación básica hasta la automatización mediante Shell Scripting.

\section{Sección I: Comandos Básicos (Navegación y Archivos)}

\subsection{Gestión de Directorios}
\textbf{Comandos: ls, cd, mkdir, pwd}
\begin{lstlisting}[language=bash]
# ls: Listar contenido
ls -lh       # Muestra archivos con tamano en formato legible (KB, MB).
ls -R        # Muestra el contenido de forma recursiva (subdirectorios).
ls -S        # Ordena los archivos por tamano.

# cd: Cambiar directorio
cd ~         # Va directamente a la carpeta personal del usuario (HOME).
cd /etc      # Cambia a la ruta absoluta de configuraciones del sistema.

# mkdir: Crear carpetas
mkdir -p proy/src/bin # Crea toda la estructura de carpetas de una vez.
\end{lstlisting}

\subsection{Manipulación de Archivos}
\textbf{Comandos: cp, mv, rm, touch, cat}
\begin{lstlisting}[language=bash]
# cp: Copiar
cp -v file1.txt backup/  # Copia y muestra que archivo se esta procesando.
cp *.jpg imagenes/       # Copia todos los archivos JPG a una carpeta.

# rm: Eliminar
rm -i archivo.txt  # Pide confirmacion antes de borrar (modo seguro).
\end{lstlisting}

\subsection{Nuevos Comandos Añadidos (Básicos)}
\begin{itemize}
    \item \textbf{alias}: Crea atajos para comandos largos.
    \begin{lstlisting}[language=bash]
    alias ll='ls -la'     # Crea el comando 'll' para ver todo detallado.
    alias ..='cd ..'      # Atajo rapido para subir de nivel.
    \end{lstlisting}
    
    \item \textbf{history}: Muestra la lista de comandos ejecutados anteriormente.
    \begin{lstlisting}[language=bash]
    history | tail -n 10  # Muestra los ultimos 10 comandos usados.
    !42                   # Ejecuta el comando numero 42 del historial.
    \end{lstlisting}

    \item \textbf{clear}: Limpia la terminal.
    \begin{lstlisting}[language=bash]
    clear                 # Limpia la pantalla visualmente.
    \end{lstlisting}
\end{itemize}

\section{Sección II: Comandos Avanzados (Sistema y Filtros)}

\subsection{Filtros de Texto}
\textbf{Comandos: grep, find, head, tail, sort, wc}
\begin{lstlisting}[language=bash]
# grep: Buscar texto
grep -i "error" log.txt   # Busca la palabra ignorando mayusculas/minusculas.
grep -v "info" log.txt    # Muestra todas las lineas que NO contienen "info".

# find: Buscar archivos
find . -type d            # Busca solo directorios en la ruta actual.
find /tmp -mtime +7       # Busca archivos modificados hace mas de 7 dias.
\end{lstlisting}

\subsection{Nuevos Comandos Añadidos (Avanzados)}
\begin{itemize}
    \item \textbf{du}: Estima el uso de espacio de archivos y directorios.
    \begin{lstlisting}[language=bash]
    du -sh * # Resumen del tamano de cada carpeta en el nivel actual.
    du -sk /var/log       # Tamano de los logs en kilobytes.
    \end{lstlisting}

    \item \textbf{df}: Informa sobre el espacio libre en el disco.
    \begin{lstlisting}[language=bash]
    df -h                 # Muestra particiones y espacio libre en formato GB/MB.
    \end{lstlisting}

    \item \textbf{free}: Muestra el estado de la memoria RAM y Swap.
    \begin{lstlisting}[language=bash]
    free -m               # Muestra el uso de memoria en Megabytes.
    \end{lstlisting}
\end{itemize}

\section{Sección III: Programación Shell (Bash Scripting)}

\subsection{Estructuras de Control}
\begin{lstlisting}[language=bash]
#!/bin/bash
# Ejemplo de bucle FOR para renombrar archivos
for f in *.txt; do
    mv "$f" "respaldo_$f"
    echo "Renombrado $f"
done
\end{lstlisting}

\subsection{Nuevos Comandos para Scripting}
\begin{itemize}
    \item \textbf{read}: Captura la entrada del usuario.
    \begin{lstlisting}[language=bash]
    read -p "Ingresa tu nombre: " nombre
    echo "Hola $nombre"
    \end{lstlisting}

    \item \textbf{sleep}: Pausa la ejecución por un tiempo determinado.
    \begin{lstlisting}[language=bash]
    sleep 5               # Espera 5 segundos antes de continuar el script.
    \end{lstlisting}

    \item \textbf{exit}: Finaliza el script con un código de estado.
    \begin{lstlisting}[language=bash]
    exit 0                # Finalizacion exitosa.
    exit 1                # Finalizacion con error.
    \end{lstlisting}
\end{itemize}

\section{Sección IV: Guía de Diagnóstico (5 Errores Comunes)}
\begin{enumerate}
    \item \textbf{Permission denied}: Falta de privilegios. Solución: \texttt{sudo [comando]} o \texttt{chmod +x [archivo]}.
    \item \textbf{Command not found}: Error tipográfico o binario fuera del PATH. Solución: Revisar \texttt{echo \$PATH}.
    \item \textbf{No such file or directory}: Ruta inexistente. Solución: Verificar con \texttt{ls} o usar rutas absolutas.
    \item \textbf{File exists}: Error al crear algo que ya está allí. Solución: Usar flags de fuerza como \texttt{mkdir -p} o \texttt{mv -f}.
    \item \textbf{Read-only file system}: El disco tiene errores y se protegió. Solución: Ejecutar \texttt{fsck} para reparar.
\end{enumerate}

\section{Sección V: Tuberías y Redireccionamientos}
Esta sección explica cómo Linux gestiona la entrada y salida de datos a través de los flujos estándar: \textbf{stdin} (0), \textbf{stdout} (1) y \textbf{stderr} (2).

\subsection{Uso de Redirecciones (\textgreater\ y \textgreater\textgreater)}
\begin{itemize}
    \item \textbf{Salida simple} (\textgreater): Sobrescribe un archivo con el resultado.
    \item \textbf{Añadir} (\textgreater\textgreater): Agrega el resultado al final del archivo sin borrar lo anterior.
\end{itemize}

\begin{lstlisting}[language=bash]
# Enviar solo los errores a un log aparte
ls /root 2> errores.log

# Guardar salida normal y errores en el mismo lugar
comando > total.txt 2>&1
\end{lstlisting}

\subsection{Uso de Tuberías (Pipes)}
Permiten conectar la salida de un comando con la entrada de otro.
\begin{lstlisting}[language=bash]
# Ejemplo: Listar procesos y buscar uno especifico
ps aux | grep "apache"

# Ejemplo: Buscar archivos grandes y ordenarlos
du -h | sort -h | tail -n 5
\end{lstlisting}
\end{document}